\section*{Hoofdstuk \ref{ch:b2:flow-balances} --- Impuls- en energiebalansen
in stromingen}

\begin{solution}{exo:b2:flow-balances:1}
$s = \qty{7.1}{cm^2}$, $v = 0.025/7.1 \times 10^{-4} = \qty{35}{m/s}$;
$F = D_mv = 25
\times 35 = \qty{880}{N}$ op het mondstuk en op de wand; bij
een terugwijkende wand:
$\rho s(v -
u)^2 = 1000 \times 7.1 \times 10^{-4} \times 25.4^2 = \qty{460}{N}$.
\end{solution}

\begin{solution}{exo:b2:flow-balances:2}
$F = 250 \times 3000 = \qty{750}{kN}$;
$a_0 = 750/60 - 9.81 = \qty{2.7}{m/s^2}$; bij \qty{20}{t}:
$37.5 - 9.8 = \qty{28}{m/s^2}$. $\Delta v = 3000\ln3 = \qty{3.3}{km/s}$ in
de vrije ruimte; min $g\tau = \qty{1.6}{km/s}$: \qty{1.7}{km/s}.
\end{solution}

\begin{solution}{exo:b2:flow-balances:3}
$S = \qty{177}{cm^2}$: $PS = \qty{7.1}{kN}$, $\rho Sv^2 = \qty{71}{N}$. Elke
component van de kracht op de knie is $PS + \rho Sv^2 = \qty{7.1}{kN}$:
\qty{10}{kN} langs de naar buiten wijzende bissectrice, waarvan $99\%$ druk.
Met open uitmonding (daar overdruk nul): $7.1$ kN langs de intreerichting en
\qty{71}{N} langs de uittreerichting: \qty{7.1}{kN}, vrijwel langs de
intree-as.
\end{solution}

\begin{solution}{exo:b2:flow-balances:4}
$F = 80 \times 550 = \qty{44}{kN}$. In vlucht:
$F = 80(550 - 240) = \qty{25}{kN}$, $Fv =
\qty{6.0}{MW}$;
$\mathcal P_{\text{kin}} = \tfrac12 \times 80 \times (550^2 - 240^2) =
\qty{9.8}{MW}$;
$\eta_p = 0.61 = 2 \times 240/790$.
\end{solution}

\begin{solution}{exo:b2:flow-balances:5}
$D_m = \qty{500}{kg/s}$. (a) $F = 1000(60 - u)$ N,
$\mathcal P = 1000u(60 - u)$; $u =
\qty{30}{m/s}$, $\omega = 37.5$ rad/s
$= \qty{360}{rpm}$. (b) \qty{900}{kW} $=
\tfrac12D_mv^2$: $100\%$. (c)
$F = D_m(v - u)(1 + \cos15^\circ) = \qty{29.5}{kN}$,
$\mathcal P
= \qty{885}{kW}$, $98\%$. (d) $\Gamma = FR = \qty{24}{kN.m}$;
bij stilstand verdubbelt de kracht ($\Gamma = \qty{48}{kN.m}$); bij $u = v$
verdwijnen kracht en koppel.
\end{solution}

\begin{solution}{exo:b2:flow-balances:6}
(a) $s = \qty{7.1}{mm^2}$ per mondstuk, elk \qty{0.10}{L/s}:
$w = \qty{14}{m/s}$. (b)
$\Gamma_0 = D_mRw = 0.2 \times 0.15 \times 14 = \qty{0.42}{N.m}$. (c)
$D_mR(w - R\omega) =
5 \times 10^{-3}$: $w - R\omega = \qty{0.17}{m/s}$,
$\omega = \qty{93}{rad/s} \approx \qty{890}{rpm}$; zonder wrijving
$R\omega = w$: \qty{900}{rpm}. (d) \qty{0.17}{m/s} achterwaarts, en daarna
nul: het water valt recht omlaag.
\end{solution}

\begin{solution}{exo:b2:flow-balances:7}
(a) $D_m = \rho Av_i$, \omterm{thm:b2:flow-balances:momentum}{impulsflux} $D_m\cdot2v_i$: $T = 2\rho Av_i^2$;
$A = \qty{78.5}{m^2}$, $T = \qty{19.6}{kN}$: $v_i = \qty{10}{m/s}$. (b)
$\mathcal P = \tfrac12D_m(2v_i)^2 = Tv_i =
\qty{200}{kW}$. (c)
$\mathcal P = T^{3/2}/\sqrt{2\rho A}$: het oppervlak verdubbelen deelt het
geïnduceerde vermogen door $\sqrt2$ --- een grotere rotor verzet meer lucht
met een lagere snelheid. (d) $A = \qty{0.20}{m^2}$,
$v_i = \sqrt{9.81/0.47} = \qty{4.6}{m/s}$, $\mathcal P
= \qty{45}{W}$:
\qty{45}{W/kg} tegen \qty{100}{W/kg} --- de schijfbelasting $T/A$ is lager.
\end{solution}

\begin{solution}{exo:b2:flow-balances:8}
(a) $v_2 = \qty{1.0}{m/s}$; Bernoulli
$\tfrac12\rho(16 - 1) = \qty{7.5}{kPa}$; werkelijk
$\rho v_2(v_1 - v_2) = \qty{3.0}{kPa}$; verlies \qty{4.5}{kPa}
$= \qty{0.46}{m}$. (b) $6$ kPa teruggewonnen, \qty{1.5}{kPa} verloren. (c)
$\tfrac12\rho v_1^2 = \qty{8}{kPa}$, \qty{0.82}{m}. (d)
$D_V = \qty{7.9}{L/s}$: \qty{35}{W} en \qty{63}{W}.
\end{solution}

\begin{solution}{exo:b2:flow-balances:9}
(a) $v = 0.01/2.83 \times 10^{-3} = \qty{3.5}{m/s}$,
$v^2/2g = \qty{0.64}{m}$; wrijving
$0.025 \times 1000 \times 0.64 = \qty{16}{m}$, hulpstukken \qty{1.3}{m}:
\qty{17.3}{m}. (b) $H = 30 + 17.3 + 0.64 = \qty{48}{m}$;
$\rho gD_VH = \qty{4.7}{kW}$; \qty{7.2}{kW} aan de as. (c) Zuigzijde (zonder
zuigverliezen): $P_0 - \rho g(5 + 0.64) = \qty{45}{kPa}$; perszijde
$45 + 470 = \qty{515}{kPa}$. (d) $P_{\text{in}} \ge \qty{2.3}{kPa}$:
$h \le (100 -
2.3)/9.81 - 0.64 = \qty{9.3}{m}$ (minder met zuigverliezen; in
de praktijk $7$--\qty{8}{m}).
\end{solution}

\begin{solution}{exo:b2:flow-balances:10}
(a) $v = \sqrt{2 \times 4 \times 10^5/1000} = \qty{28}{m/s}$;
$s = \qty{3.8}{cm^2}$, $D_m =
\qty{10.8}{kg/s}$,
$F = D_mv = 2\Delta P\,s = \qty{300}{N}$;
$a = 300/1.1 - 9.8 =
\qty{260}{m/s^2}$. (b) Lucht op \qty{1.5}{L}:
$P = \qty{3.3}{bar}$,
$v = \sqrt{2 \times 2.3 \times
10^5/1000} = \qty{22}{m/s}$. (c) Gemiddelde
uitstroomsnelheid $\approx \qty{22}{m/s}$: duur
$\approx 10^{-3}
/(3.8 \times 10^{-4} \times 22) = \qty{0.12}{s}$;
\omterm{prop:b2:flow-balances:pelton}{raketvergelijking} met $v_e \approx \qty{22}{m/s}$:
$\Delta v \approx 22\ln(1.1/0.1) = \qty{50}{m/s}$. (d) De resterende
\qty{2.5}{bar} lucht blaast naar buiten en voegt er nog wat aan toe; te veel
water laat weinig lucht over en de druk stort vroeg in, te weinig water laat
weinig massa over om weg te werpen --- ongeveer een derde van het volume is
het beste.
\end{solution}

\begin{solution}{exo:b2:flow-balances:11}
(a)
$\Gamma = D_mr_1v_{\theta1} = 3 \times 10^4 \times 20 = \qty{600}{kN.m}$,
$\omega = 15.7$ rad/s, $\mathcal P = \qty{9.4}{MW}$. (b)
$H = \mathcal P/\rho gD_V = \qty{32}{m}$. (c)
$\Gamma = 3 \times
10^4(20 - 1.5) = \qty{555}{kN.m}$,
$\mathcal P = \qty{11.6}{MW}$ (drukhoogte \qty{39}{m}); de uittreedraaiing
voert $v_{\theta2}^2/2g = \qty{0.46}{m}$ mee, ongeveer $1\%$. (d) Een straal
van Pelton staat op atmosferische druk: zij heeft een grote vervalhoogte
nodig om snel te zijn, en haar debiet wordt door het mondstuk begrensd; een
loopwiel van Francis is gevuld met water onder druk en verwerkt grote
debieten bij matige vervalhoogten.
\end{solution}

\begin{solution}{exo:b2:flow-balances:12}
(a) $v_1h_1 = v_2h_2 = q$. (b) Hydrostatische kracht per eenheid breedte op
een doorsnede $\int_0^h\rho g(h - z)\dd z = \tfrac12\rho gh^2$; impuls:
$\tfrac12\rho g(h_1^2
- h_2^2) = \rho q(v_2 - v_1)$; deel door $\rho$ en
gebruik $q = v_1h_1 = v_2h_2$. (c) Met $v_2 = v_1h_1/h_2$ en $h_1 \ne h_2$:
$\tfrac12g(h_1 + h_2)h_2 = v_1^2h_1$, d.w.z.\
$(h_2/h_1)^2 + h_2/h_1 - 2\mathrm{Fr}_1^2 = 0$; $h_2 > h_1$ vereist
$\mathrm{Fr}_1 > 1$. (d)
$\Delta H = (h_1 + v_1^2/2g) - (h_2 + v_2^2/2g) = (h_2 - h_1)^3/4h_1h_2$.
$\mathrm{Fr}_1 = 2.86$, $h_2/h_1 = 3.57$, $h_2 = \qty{0.71}{m}$,
$v_2 = \qty{1.1}{m/s}$, $\Delta H = 0.514^3/0.57 = \qty{0.24}{m}$;
$\rho gq\,\Delta H = \qty{1.9}{kW}$ per meter breedte.
\end{solution}

\begin{solution}{pb:b2:flow-balances:1}
\textbf{1.} $F = (D_a + D_f)v_e = 51 \times 600 = \qty{30.6}{kN}$.

\textbf{2.} $F = 51 \times 600 - 50 \times 250 = \qty{18.1}{kN}$;
$Fv = \qty{4.5}{MW}$.

\textbf{3.}
$\mathcal P_{\text{kin}} = \tfrac12 \times 51 \times 600^2 - \tfrac12 \times 50 \times 250^2 =
\qty{7.6}{MW}$;
$\eta_p = 0.59$. Met $D_f = 0$:
$\eta_p = 2v(v_e - v)/(v_e^2 - v^2)
= 2v/(v_e + v) = 500/850$.

\textbf{4.} Warmte \qty{43}{MW}: thermisch $7.6/43 = 18\%$; totaal $4.5/43 = 10\%$.

\textbf{5.} $\tfrac12 \times 500(v_e'^2 - 250^2) = 7.6 \times 10^6$:
$v_e' = \qty{305}{m/s}$; $F = 500 \times 55 = \qty{27.5}{kN}$;
$\eta_p = 500/555 = 0.90$. Meer lucht met een lagere snelheid verzetten
geeft bij hetzelfde vermogen meer stuwkracht: vandaar de fans.

\textbf{6.} Zij heeft de buitenlucht zowel als werkmedium als als oxidator
nodig; een raket draagt beide zelf mee.

\textbf{7.} De lucht komt binnen met \qty{70}{m/s} (achterwaarts in het
stelsel van de motor) en vertrekt met een voorwaartse component
$150\cos45^\circ = \qty{106}{m/s}$: de impuls die aan de lucht wordt gegeven
is voorwaarts, $500(106 + 70) = \qty{88}{kN}$ remkracht; dezelfde stroom
achterwaarts uitgestoten zou $500(150 - 70) =
\qty{40}{kN}$ stuwkracht
geven.

\textbf{8.} Over $\dd t$ geldt voor het systeem (raket $m$ + gas $D_m\dd t$
uitgestoten met $v - v_e$):
$(m - D_m\dd t)(v + \dd v) + D_m\dd t(v - v_e) - mv = -mg\,\dd t$, en dus
$m\,\dd v/\dd t = D_mv_e - mg$.

\textbf{9.} $D_m = \qty{1333}{kg/s}$; $F = \qty{4.0}{MN}$;
$a = 13.3 - 9.8 =
\qty{3.5}{m/s^2}$ bij het opstijgen,
$40 - 9.8 = \qty{30}{m/s^2}$ bij het uitbranden.

\textbf{10.} $v = v_e\ln(m_0/m) - gt$:
$3000\ln3 - 9.81 \times 150 = 3296 - 1472 =
\qty{1.8}{km/s}$; het
zwaartekrachtverlies bedraagt \qty{1.5}{km/s}.

\textbf{11.} Trap 1: $3000\ln(300/200) = \qty{1.2}{km/s}$; werp \qty{20}{t}
af; trap 2: $3000\ln(180/80) = \qty{2.4}{km/s}$: \qty{3.6}{km/s} tegen
\qty{3.3}{km/s} voor één trap (met zwaartekracht wordt het verschil groter).

\textbf{12.} $v_e = 7800/\ln3 = \qty{7.1}{km/s}$: geen enkele chemische
stuwstof haalt dat (waterstof--zuurstof geeft \qty{4.5}{km/s}); trappen zijn
onvermijdelijk.

\textbf{13.} De verticale \omterm{thm:b2:flow-balances:momentum}{impulsflux} van de straal, $D_mv = \qty{4.0}{MN}$,
wordt weggenomen en er wordt een even grote horizontale gemaakt: een
neerwaartse kracht van \qty{4}{MN} plus \qty{4}{MN} zijwaarts, samen
\qty{5.7}{MN}.

\textbf{14.} De balans op de motor bevat $(P_e - P_0)S_e$ op de
uitlaatdoorsnede; omdat $P_0$ met de hoogte daalt, groeit die term: de
stuwkracht neemt in vacuüm met zo'n $10$--$15\%$ toe.

\textbf{15.} $D_m = \qty{3e4}{kg/s}$ per wiel:
$F = D_m(v - u)(1 + \cos15^\circ) =
1.97 \times 3 \times 10^4(62.6 - u)$;
$\mathcal P = Fu$, maximaal bij $u = v/2 = \qty{31.3}{m/s}$.

\textbf{16.} $\omega = u/R = 20.9$ rad/s $= \qty{199}{rpm}$;
$F = \qty{1.85}{MN}$, $\Gamma = FR
= \qty{2.8}{MN.m}$.

\textbf{17.} $f = pn/60$: $p = 3000/199 \approx 15$ poolparen, $n = \qty{200}{rpm}$.

\textbf{18.} $\mathcal P = Fu = \qty{58}{MW}$; vermogen van de straal
$\tfrac12D_mv^2 = \qty{59}{MW}$: $98\%$;
$2 \times 58 \times 0.95 = \qty{110}{MW}$.

\textbf{19.} Intree: $Rv = \qty{94}{m^2/s}$ per kg. Uittree: absolute
azimutale snelheid $u - (v - u)\cos15^\circ = \qty{1.1}{m/s}$, d.w.z.\
\qty{1.6}{m^2/s}. $\Gamma = D_m(94 -
1.6) = \qty{2.8}{MN.m}$ --- zoals in
16.

\textbf{20.} Naarmate $u \to v$ verdwijnen de relatieve snelheid, de kracht
en het koppel: niets remt het wiel meer af, dat op hol slaat; de deflector
werpt de straal binnen een seconde van de schoepen af, terwijl de
naaldafsluiter langzaam genoeg sluit om waterslag te vermijden.

\textbf{21.} $v = \qty{5.7}{m/s}$, $v^2/2g = \qty{1.6}{m}$; verlies
$0.015 \times 133 \times 1.6 =
\qty{3.3}{m}$; opvoerhoogte
$200 + 3.3 + 1.6 = \qty{205}{m}$.

\textbf{22.} $\rho gD_VH = 9810 \times 40 \times 205 = \qty{80}{MW}$;
\qty{91}{MW} elektrisch.

\textbf{23.}
$P = P_0 + \rho g(200 + 3.3) + \tfrac12\rho v^2 \approx \qty{21}{bar}$
absoluut, tegen \qty{20.3}{bar} in turbinebedrijf: de verliezen tellen nu
bij de statische drukhoogte op in plaats van eraf te gaan.

\textbf{24.} $V = 40 \times 28800 = \qty{1.15e6}{m^3}$,
$E = \rho gV \times 200 = \qty{2.3e12}{J}
= \qty{630}{MWh}$; verbruikt
$91 \times 8 = \qty{730}{MWh}$; teruggewonnen
$630 \times 0.9 \times
0.95 = \qty{540}{MWh}$: heen-en-terugrendement
$74\%$.

\textbf{25.} Impuls: de stuwkracht van de straalmotor en van de raket, de
kracht op het peltonwiel. Impulsmoment: het turbinekoppel (Euler). Energie:
de opvoerhoogte en het vermogen van de pomp, het rendement van de opslag.
Massa: elk debiet, en het massaverlies van de raket.
\end{solution}
